\chapter{Engineering Decisions and Limitations}
\label{ch:decisions}

\section{Why this architecture was selected}
\begin{description}
 \item[Use the supplied GP7 CAD-derived description.] This preserves the
 assignment's industrial-robot model and lets the group focus on kinematics,
 control, and task execution. The authoritative task URDF is explicit and
 inspectable; separate CAD URDF exports remain archival inputs.
 \item[Derive standard DH, execute URDF FK.] A conventional DH table supports
 lecture-level frame and FK explanation. The runtime reads exact origins and
 axes from URDF so it does not duplicate dimensions in control code. A
 separate verifier independently checks the table against the URDF chain.
 \item[Use bounded numerical IK.] SciPy's least-squares solver is concise and
 consumes limits parsed from the model. Warm-starting supports waypoint
 continuity, while restart seeds help find the first feasible approach. The
 trade-off is no closed-form branch analysis and no collision constraints.
 \item[Use quintic interpolation.] The scalar profile has rest-to-rest
 endpoint velocity and acceleration and is straightforward to sample into a
 ROS trajectory. The project applies a sampled velocity cap and stretches
 time, but does not establish acceleration/jerk limits.
 \item[Use ROS 2 control actions.] The controller manager, broadcaster, and
 joint trajectory controller separate command generation from simulation
 actuation. The task can await goal results, inspect feedback, and cancel.
 \item[Use a simple adaptive gripper.] One commanded joint with mimic metadata
 is simpler than six independent actuators and visually resembles the
 supplied end-effector CAD. The selected Gazebo engine's unsupported mimic
 constraint makes that design unverified dynamically.
 \item[Use detachable-joint retention for task completion.] It guarantees a
 repeatable simulated carry and lets the group validate sequencing and pose
 bookkeeping. It bypasses the physical grasp mechanics, so it must be labeled
 idealized rather than interpreted as force-closure performance.
\end{description}

\section{Reachability trade-off}
The assignment sheet's Box A at $(1,1,0)$ m is outside the described reach.
The simulation keeps a fixed reachable Box A at $(0.55,0.35,0)$ only under the
lecturer's later verbal clarification. This discrepancy must be stated in
presentations and submissions. Do not confuse the clarification with a
written assignment amendment. Do not describe seeded Box B test coordinates
as moving Box A.

The far-mode maximum radius is an optimization subject to $z=0.025$ and fixed
downward orientation, not an unrestricted reach sphere. Its exact boundary
is a straight-elbow singularity. A 5 mm inward margin and complete route
preflight reduce this modeled edge case, but do not quantify real-world
tracking margin.

\section{Numerical and modeling limits}
\begin{itemize}
 \item The IK residual mixes metres and radians with no unit weighting; solver
 convergence may depend on scale and seed.
 \item The solver has finite evaluations and can fail, especially at
 singularities or on infeasible targets. Multiple restarts do not prove
 completeness.
 \item Path continuity uses a maximum adjacent joint displacement threshold;
 no Jacobian conditioning check is computed.
 \item The selected Cartesian path is straight in tool position with fixed
 orientation. It is not a global collision-aware path.
 \item The only trajectory bound beyond URDF speed is a sampled 0.5 rad/s
 cap. Acceleration and jerk are neither bounded by code nor measured in the
 validation evidence.
 \item URDF mass and inertia values are repeated estimates. Mesh/collision
 accuracy and CAD-to-URDF equivalence are not established.
 \item Gripper mimic behavior is unsupported by the reported physics setup;
 the fixed-joint plugin creates idealized object retention.
 \item Pose feedback is sampled at 50 Hz. The reported orientation bound is
 the maximum over received samples, not proof over every physics step.
\end{itemize}

\section{What a more physical system would require}
An engineering route beyond the current proof-of-concept would include
manufacturer-verified link inertias and calibration; collision geometry
validation; support for mimic joints or independent finger control; contact
friction and force models; grasp-force/current/sensor feedback; robust
collision checking; joint trajectory logging; and systematic repeated runs
with disturbances. A physical deployment additionally requires safety-rated
limits, real hardware interfaces, emergency stop, workspace risk review, and
commissioning; the project's \code{/gp7/stop} topic is not sufficient.

\section{Implemented, inspected, and future categories}
\begin{table}[htbp]
\centering\small
\caption{Evidence labels to use when explaining decisions.}\label{tab:evidence-levels}
\begin{tabularx}{\textwidth}{p{.23\textwidth}X X}
\toprule Category & Examples & Accurate wording\\\midrule
Implemented and validated in recorded simulation & Full sequence; controller activation; box pose checking; plugin attach/detach; successful placement checks & ``The Gazebo task completed under the recorded setup.''\\
Implemented and offline-checked & Bounded URDF IK; trajectory samples; speed and joint-limit tests; gripper description transformation & ``The code/test checks this path offline.''\\
Mathematical explanation & Jacobian concepts, singularity theory, rigid transforms, quintic derivatives & ``This theory explains behavior; no explicit Jacobian controller is coded.''\\
Future enhancement & Contact-force grasp, real emergency stop, collision planner, physical accuracy & ``Not implemented or validated in this repository.''\\
\bottomrule
\end{tabularx}
\end{table}
